% Extracción quirúrgica del propietario V2 05c: líneas 99--335.
% El resto permanece en este archivo como procedencia, pero no se publica.
\iffalse
\chapter[Internal Gauge Macrostructure]{Internal Gauge Macrostructure Generated by the Discrete Continuum}
\label{chap:gauint-macroestructura}

\begin{thesisbox}
The macrostructure of this chapter is constructed solely from a
surviving restriction of \(\mathcal C_{\rm cont}^{\rm disc}\). The numbers
\(4\), \(6\), and \(8\), the screen, the four involutions, the
projective measure, the transfer operator, the incidence mode, and the cellular
terminal are derived within that restriction. No subsequent evaluation
forms part of the domain or of the arrows that follow.
\end{thesisbox}

\section{The projective seed and the internal derivation of \texorpdfstring{\(4/6/8\)}{4/6/8}}

Let \(V_{\rm TPK}=\mathbf F_3^2\), with the ordered basis inherited from the two
sheets of the TPK state. We define the set of directions
\begin{equation}\label{eq:gauint-cuatro-direcciones}
 I_{\rm gau}:=\mathbf P^1(\mathbf F_3)
 =\{L_\infty,L_0,L_1,L_{-1}\},
\end{equation}
where
\[
 L_\infty=[1:0],\qquad L_0=[0:1],\qquad
 L_1=[1:1],\qquad L_{-1}=[1:-1].
\]
Therefore,
\begin{equation}\label{eq:gauint-cardinal-cuatro}
 |I_{\rm gau}|=\frac{3^2-1}{3-1}=4.
\end{equation}
A set of four indices is not introduced: it appears as the quotient of
the eight nonzero vectors by the action of \(\mathbf F_3^\times\).

The unoriented incidences are
\begin{equation}\label{eq:gauint-seis-incidencias}
 E_{\rm gau}:=\binom{I_{\rm gau}}2,
 \qquad |E_{\rm gau}|=\binom42=6,
 \qquad Q_{\rm gau}:=\mathbf F_3^{E_{\rm gau}}.
\end{equation}
The oriented screen duplicates each direction, not each incidence:
\begin{equation}\label{eq:gauint-ocho-caras}
 \mathscr B_{\rm gau}:=I_{\rm gau}\times\{+,-\},
 \qquad |\mathscr B_{\rm gau}|=2|I_{\rm gau}|=8.
\end{equation}
Thus, \(4\), \(6\), and \(8\) arise, respectively, from directions, pairs of
directions, and oriented faces of the same seed.

For each \(L\in I_{\rm gau}\), let \(\vartheta_L\) be the screen
involution that exchanges \((L,+)\) and \((L,-)\) and fixes the other six
faces. Then
\begin{equation}\label{eq:gauint-cuatro-involuciones}
 \vartheta_L^2=I,
 \qquad
 \vartheta_L\vartheta_{L'}=\vartheta_{L'}\vartheta_L
 \quad(L\ne L').
\end{equation}
The four involutions have been obtained prior to defining measure or operator.

\section{Incidence, internal mode and non-constant witness}

The incidence application is
\begin{equation}\label{eq:gauint-incidencia-b}
 B:\mathbf F_3^{I_{\rm gau}}\longrightarrow Q_{\rm gau},
 \qquad
 (Ba)_{\{L,L'\}}:=a_L+a_{L'}.
\end{equation}
If
\[
 s(q):=\sum_{e\in E_{\rm gau}}q_e\in\mathbf F_3,
\]
each direction appears exactly three times in the sum of the six
incidences, and therefore
\begin{equation}\label{eq:gauint-incidencia-invariante}
 s(Ba)=3\sum_{L\in I_{\rm gau}}a_L=0,
 \qquad
 s(q+Ba)=s(q).
\end{equation}

On \(Q_{\rm gau}\), only the normalized counting average is used.
The incidence mode is the rational function
\begin{equation}\label{eq:gauint-wc-definicion}
 W_c(q):=\mathbf 1_{\{s(q)=0\}}-\frac13.
\end{equation}
Since \(s:Q_{\rm gau}\to\mathbf F_3\) is surjective, each fiber has \(3^5\) elements.
Exact counting then yields,
\begin{equation}\label{eq:gauint-wc-momentos}
 \mathbb E_{Q}W_c=0,
 \qquad
 \mathbb E_{Q}W_c^2=\frac29,
 \qquad
 \mathbb E_{Q}W_c^3=\frac{2}{27}.
\end{equation}
Furthermore, \eqref{eq:gauint-incidencia-invariante} gives
\begin{equation}\label{eq:gauint-wc-invariancia}
 W_c(q+Ba)=W_c(q).
\end{equation}
Therefore, \(W_c\) is nonzero, centered, and invariant under the internally
generated incidence.

\fi
\section{Incidence, carry curvature, and common domain of naturalness}

For \(x\in\mathcal C_{\rm cont}^{\rm disc}\), let us denote by
\(Y_m(x)\) the finite set of gauge prefixes of depth \(m\) that
simultaneously conserve
\begin{equation*}
\begin{gathered}
 \text{fiber, extension relations, numbers, orientation, carry, memory,}\\
 \text{boundary, route, multisection, survival, terminal and reader}.
\end{gathered}
\end{equation*}
There exists a collar neighborhood application.
\begin{equation}\label{eq:gauint-collar}
 q_m:Y_m(x)\longrightarrow Q_{\rm gau}
\end{equation}
which records the six incidences without erasing the history that produces them.
For \(\ell\ge m\), the truncation is written
\[
 \tau_{\ell,m}:Y_\ell(x)\longrightarrow Y_m(x).
\]

For \(y,z\in Y_m(x)\), let
\begin{equation}\label{eq:gauint-todas-extensiones}
 \operatorname{Ext}^{\rm adm}_{m}(y,z)
\end{equation}
the set of \emph{all} admissible TPK extensions from \(y\) to \(z\).
An extension belongs to this set only if it transports the entire fiber,
the relations \(\operatorname{Ext}_\Gamma\), the native and operatorial carry,
the orientation, the accumulated memory, the boundary, the route, and the
survival condition. We define
\begin{equation}\label{eq:gauint-conteo-kernel}
 a_m(y,z):=\#\operatorname{Ext}^{\rm adm}_m(y,z),
 \qquad
 d_m(y):=\sum_{z\in Y_m(x)}a_m(y,z),
\end{equation}
and, when \(d_m(y)>0\),
\begin{equation}\label{eq:gauint-kernel-normalizado}
 k_m(y,z):=\frac{a_m(y,z)}{d_m(y)}.
\end{equation}
Thus, \(k_m\) is a normalized count of complete extensions, not a
function introduced after projecting the state.

The restriction \(\mathcal C_{\rm cont}^{\rm gau}\) is formed by the
surviving histories \(x\) for which the following conditions
hold at every depth.
\begin{enumerate}
 \item The sets \(Y_m(x)\) are nonempty, \(d_m(y)>0\), and the
 truncations \(\tau_{\ell,m}\) are surjective.
 \item There exists a nonempty multisection of strictly positive rational
 weights \(\mu_m:Y_m(x)\to\mathbf Q_{>0}\) such that
 \begin{equation}\label{eq:gauint-medida-proyectiva}
  \sum_{y\in Y_m}\mu_m(y)=1,
  \qquad
  \mu_m(y)=\sum_{\tau_{\ell,m}\widetilde y=y}\mu_\ell(\widetilde y).
 \end{equation}
 \item Each \(\mu_m\) is stationary:
 \begin{equation}\label{eq:gauint-estacionariedad}
  \sum_{y\in Y_m}\mu_m(y)k_m(y,z)=\mu_m(z).
 \end{equation}
 \item The marginal of \((q_m)_*\mu_m\) is the uniform counting measure on
 \(Q_{\rm gau}\). This condition makes the calculations
 of \eqref{eq:pdfv2-cor-wc-momentos} literally applicable to the lift \(W_c\circ q_m\).
 \item Direct lumpability holds
 \begin{equation}\label{eq:gauint-lump-directa}
  \sum_{\tau_{\ell,m}\widetilde z=z}
   k_\ell(\widetilde y,\widetilde z)
  =k_m(\tau_{\ell,m}\widetilde y,z)
 \end{equation}
 for every \(\widetilde y\in Y_\ell\) and \(z\in Y_m\).
 \item For the adjoint kernel
 \begin{equation}\label{eq:gauint-kernel-adjunto}
  k_m^*(z,y):=\frac{\mu_m(y)k_m(y,z)}{\mu_m(z)},
 \end{equation}
 the adjoint lumpability condition is satisfied
 \begin{equation}\label{eq:gauint-lump-adjunta}
  \sum_{\tau_{\ell,m}\widetilde y=y}
   k_\ell^*(\widetilde z,\widetilde y)
  =k_m^*(\tau_{\ell,m}\widetilde z,y).
 \end{equation}
 \item The categories of prefixes possess the TPK root as an initial object,
 so that their nerve has the explicit terminal contraction of
 \eqref{eq:gauint-sdr} below.
\end{enumerate}
These conditions define the typed domain; outside it, a measure or natural
operator is not published by mere declaration.

\section{Projective measure, kernel, and direct and adjoint naturality}

Having fixed a branch of the multisection \(\{\mu_m\}_m\), let
\begin{equation}\label{eq:gauint-hm}
 H_m:=\operatorname{Fun}(Y_m,\mathbf Q),
 \qquad
 \langle f,g\rangle_m:=\sum_{y\in Y_m}\mu_m(y)f(y)g(y).
\end{equation}
The operator generated by all extensions is
\begin{equation}\label{eq:gauint-km}
 (K_mf)(y):=\sum_{z\in Y_m}k_m(y,z)f(z).
\end{equation}
The Markov equations and stationarity give
\begin{equation}\label{eq:gauint-markov-estacionario}
 K_m\mathbf1=\mathbf1,
 \qquad
 K_m^*\mathbf1=\mathbf1,
 \qquad
 \|K_mf\|_m\le\|f\|_m.
\end{equation}

The lift and the truncated expectation are
\begin{equation}\label{eq:gauint-lift-esperanza}
 J_{\ell,m}f:=f\circ\tau_{\ell,m},
 \qquad
 E_{m,\ell}:=J_{\ell,m}^*.
\end{equation}
The projectivity of \(\mu\) implies
\begin{equation}\label{eq:gauint-ej-identidad}
 E_{m,\ell}J_{\ell,m}=I_{H_m}.
\end{equation}
The two lumpability conditions discharge, respectively,
\begin{equation}\label{eq:gauint-naturalidad-directa-adjunta}
 K_\ell J_{\ell,m}=J_{\ell,m}K_m,
 \qquad
 K_\ell^*J_{\ell,m}=J_{\ell,m}K_m^*.
\end{equation}
The second equality is not inferred from the first without the attached condition:
both remain registered in the domain.

The positive internal operator is defined
\begin{equation}\label{eq:gauint-transferencia-t}
 T_m:=K_m^*K_m.
\end{equation}
Then
\begin{equation}\label{eq:gauint-t-positivo-natural}
 \langle f,T_mf\rangle_m=\|K_mf\|_m^2\ge0,
 \qquad
 T_\ell J_{\ell,m}=J_{\ell,m}T_m.
\end{equation}

\section{Common law of \(8\) faces and \(4\) Fubini squares}

On \(Y_m^{\mathscr B_{\rm gau}}\), we define the rational mass
\begin{equation}\label{eq:gauint-ley-ocho-caras}
 \Omega_m^{(8)}(\mathbf y)
 :=\sum_{z\in Y_m}\mu_m(z)
 \prod_{(L,\varepsilon)\in\mathscr B_{\rm gau}}
 k_m(z,y_{L,\varepsilon}).
\end{equation}
It is a probability because each factor is Markov. All faces arise
from the same latent state \(z\) and the same kernel; they are not eight independently
added measures.

For functions \(f_{L,+},f_{L,-}\in H_m\), finite Fubini gives
\begin{equation}\label{eq:gauint-fubini-ocho}
\begin{aligned}
 &\sum_{\mathbf y}\Omega_m^{(8)}(\mathbf y)
  \prod_{L\in I_{\rm gau}}
  f_{L,+}(y_{L,+})f_{L,-}(y_{L,-})\\
 &\qquad=
 \sum_{z\in Y_m}\mu_m(z)
 \prod_{L\in I_{\rm gau}}
 (K_mf_{L,+})(z)(K_mf_{L,-})(z).
\end{aligned}
\end{equation}
Upon taking \(f_{L,+}=f_{L,-}=f_L\), the four squares appear
simultaneously
\begin{equation}\label{eq:gauint-cuatro-cuadrados}
 \sum_{z\in Y_m}\mu_m(z)
 \prod_{L\in I_{\rm gau}}(K_mf_L(z))^2.
\end{equation}
The marginal of any opposite pair satisfies
\begin{equation}\label{eq:gauint-marginal-t}
 \sum_{y_+,y_-}\Omega_{m,L}^{(2)}(y_+,y_-)
 f(y_+)g(y_-)
 =\langle K_mf,K_mg\rangle_m
 =\langle f,T_mg\rangle_m.
\end{equation}
The commutativity of the product in \eqref{eq:gauint-ley-ocho-caras} proves
\begin{equation}\label{eq:gauint-involuciones-ley}
 (\vartheta_L)_*\Omega_m^{(8)}=\Omega_m^{(8)}
 \qquad(L\in I_{\rm gau}).
\end{equation}

\section{Exact internal operator and self-test witness}

The projection onto the constant sector and its complement are
\begin{equation}\label{eq:gauint-pq}
 P_{\Omega,m}f:=\langle\mathbf1,f\rangle_m\mathbf1,
 \qquad
 Q_m:=I-P_{\Omega,m}.
\end{equation}
The internal cell operator is defined by
\begin{equation}\label{eq:gauint-d-hmt}
 D_m^{\rm HMT}:=3Q_m.
\end{equation}
The coefficient \(3\) is the cardinality of the residue field used in
\eqref{eq:pdfv2-cor-proyectores-cuatro-direcciones}; it does not represent an
added scale.
For the lift
\[
 W_{c,m}:=W_c\circ q_m
\]
one has, by \eqref{eq:pdfv2-cor-wc-momentos},
\begin{equation}\label{eq:gauint-w-eigen}
 P_{\Omega,m}W_{c,m}=0,
 \qquad
 D_m^{\rm HMT}W_{c,m}=3W_{c,m},
 \qquad
 \|W_{c,m}\|_m^2=\frac29.
\end{equation}
Thus, the complementary sector possesses a nonzero witness constructed by the
incidence of the six edges.

\section{Cellular terminal and strong retractor}

Let \(C_*(N\mathsf Y_m;A_n)\) be the normalized chain complex of the nerve of the
prefix category, and let \(y_*\) be its initial root. The inclusion and the
augmentation are
\[
 \iota_m:A_n[0]\longrightarrow C_*(N\mathsf Y_m;A_n),
 \quad \iota_m(1)=[y_*],
 \qquad
 \epsilon_m:C_*(N\mathsf Y_m;A_n)\longrightarrow A_n[0].
\]
The initial object provides the extra degeneracy
\begin{equation}\label{eq:gauint-homotopia-celular}
 H_m[y_0\to\cdots\to y_q]
 :=[y_*\to y_0\to\cdots\to y_q],
\end{equation}
with the normalized convention when a degeneration appears. It is verified
\begin{equation}\label{eq:gauint-sdr}
 \epsilon_m\iota_m=I,
 \qquad
 \partial H_m+H_m\partial=I-\iota_m\epsilon_m,
 \qquad
 H_m^2=H_m\iota_m=\epsilon_mH_m=0.
\end{equation}
This cellular terminal is distinct from the constant projection
\(P_{\Omega,m}\): the former contracts the nerve, while the latter separates the
sectors of the function space. Both are preserved.

\iffalse
\section{The thirteen coordinates of the gauge constraint}

For each \((m,n)\), it is defined
\begin{equation}\label{eq:gauint-d-trece}
 D_{{\rm gau},m,n}(x):=
 (\mathcal F,\operatorname{Ext}_\Gamma,\operatorname{Rel},
 \mathcal N,\operatorname{or},\operatorname{Carr},\operatorname{Mem},
 \partial,\gamma,\operatorname{Msect},\operatorname{Surv},
 \operatorname{Term},\mathcal L)_{{\rm gau},m,n}(x),
\end{equation}
with the following content.
\begin{enumerate}
 \item \(\mathcal F=(V_{\rm TPK},I_{\rm gau},E_{\rm gau},
 Q_{\rm gau},\mathscr B_{\rm gau},Y_m,H_m)\).
 \item \(\operatorname{Ext}_\Gamma\) contains all the sets
 \(\operatorname{Ext}^{\rm adm}_m(y,z)\), truncations, lifts,
 expectations, and their compositions.
 \item \(\operatorname{Rel}\) contains \(B\), \(sB=0\), the four
 involutions, Markov, stationarity, and both lumpabilities.
 \item \(\mathcal N=(3,4,6,8,m,n,|Y_m|,a_m,d_m)\); those numbers do not
 replace any of the remaining coordinates.
 \item \(\operatorname{or}\) preserves the faces \(+/-\), the four
 inversions \(\vartheta_L\), and the reversal of each path.
 \item \(\operatorname{Carr}\) contains the complete carry of every extension counted in \(a_m(y,z)\), including the carry already propagated.
 \item \(\operatorname{Mem}\) conserves the full prefix, the common
 state \(z\) of \(\Omega_m^{(8)}\), and the memory of its eight extensions.
 \item \(\partial=(q_m,\partial_{N\mathsf Y_m},Q_m,H_m)\) preserves
 collar boundary, cellular boundary, and complement.
 \item \(\gamma\) is the oriented simplex of the nerve that records the route
 and not only its endpoints.
 \item \(\operatorname{Msect}\) is the family of all compatible
 projective stationary measures and of all the extensions that
 enter into the counts; no retrospective branch is selected.
 \item \(\operatorname{Surv}\) records nonemptiness, arbitrary
 depth, projectivity, both lumpabilities, and the persistence of the
 uniform collar marginal.
 \item \(\operatorname{Term}=(\iota_m,\epsilon_m,H_m;
 P_{\Omega,m},Q_m)\) keeps the two terminals separate.
 \item \(\mathcal L=(\Omega_m^{(8)},K_m,T_m,D_m^{\rm HMT},W_{c,m})\)
 is the internal reader posterior to the construction of the history.
\end{enumerate}

For \(\ell\ge m\) and \(n'\ge n\), all the coordinates satisfy
\begin{equation}\label{eq:gauint-naturalidad-trece}
 u^{\rm gau}_{(\ell,n'),(m,n)}D_{{\rm gau},\ell,n'}(x')
 =D_{{\rm gau},m,n}(\tau_{\ell,m}x').
\end{equation}

\section{Assembler, publisher, and non-ablative chain}

Let
\begin{equation}\label{eq:gauint-graph-d}
 \mathcal G_{\rm gau}:=\operatorname{Graph}(D_{\rm gau}),
 \qquad
 \operatorname{Res}_{\rm gau}(x):=(x,D_{\rm gau}x).
\end{equation}
The dependent multisection \(\chi_{\rm gau}(x)\) contains all compatible branches
of collar, measure, and extension. We define
\begin{equation}\label{eq:gauint-x-dependiente}
 X_{\chi,{\rm gau}}
 :=\{(\chi_{\rm gau}(x),x,D_{\rm gau}x):
 x\in\mathcal C_{\rm cont}^{\rm gau}\},
\end{equation}
\begin{equation}\label{eq:gauint-prx}
 \operatorname{pr}_{X,{\rm gau}}(x,D_{\rm gau}x)
 :=(\chi_{\rm gau}(x),x,D_{\rm gau}x).
\end{equation}

The raw assembler
\begin{equation}\label{eq:gauint-ensamblador-crudo}
 a_{\rm gau}:X_{\chi,{\rm gau}}\longrightarrow
 \mathscr M_{\rm gau}^{\rm amb}
\end{equation}
produce the complete family
\[
 \bigl(I,E,\mathscr B,B,W_c;
 Y_m,\mu_m,k_m,J_{\ell,m},E_{m,\ell};
 K_m,T_m,\Omega_m^{(8)},P_{\Omega,m},Q_m,D_m^{\rm HMT};
 \iota_m,\epsilon_m,H_m\bigr)_{m,n}.
\]
The macroobject preserves its entry through
\begin{equation}\label{eq:gauint-graph-a}
 \mathfrak M_{\rm gau}:=\operatorname{Graph}(a_{\rm gau}),
 \qquad
 \mathcal A_{\rm gau}(z):=(z,a_{\rm gau}z).
\end{equation}

The raw publisher
\begin{equation}\label{eq:gauint-publicador-crudo}
 p_{\rm gau}:\mathfrak M_{\rm gau}\longrightarrow
 \mathscr Y_{\rm gau}
\end{equation}
publishes jointly the identities
\[
 \begin{gathered}
 |I|=4,quad |E|=6,quad |\mathscr B|=8,quad sB=0,
 \quad W_c(q+Ba)=W_c(q),\\
 K_\ell J=JK_m,quad K_\ell^*J=JK_m^*,
 \quad T_m=K_m^*K_m,\\
 (\vartheta_L)_*\Omega_m^{(8)}=\Omega_m^{(8)},
 \quad \text{the four quadratic Fubini factorizations},\\
 D_m^{\rm HMT}=3Q_m,quad D_m^{\rm HMT}W_{c,m}=3W_{c,m},
 \quad \partial H_m+H_m\partial=I-\iota_m\epsilon_m.
 \end{gathered}
\]
Finally,
\begin{equation}\label{eq:gauint-graph-p}
 \mathcal C_{\rm gau}^{\rm HMT}:=\operatorname{Graph}(p_{\rm gau}),
 \qquad
 \mathcal P_{\rm gau}(M):=(M,p_{\rm gau}M).
\end{equation}
The complete chain is
\begin{equation}\label{eq:gauint-cadena-completa}
 \mathcal C_{\rm cont}^{\rm disc}\supset
 \mathcal C_{\rm cont}^{\rm gau}
 \xrightarrow{\operatorname{Res}_{\rm gau}}\mathcal G_{\rm gau}
 \xrightarrow{\operatorname{pr}_{X,{\rm gau}}}X_{\chi,{\rm gau}}
 \xrightarrow{\mathcal A_{\rm gau}}\mathfrak M_{\rm gau}
 \xrightarrow{\mathcal P_{\rm gau}}\mathcal C_{\rm gau}^{\rm HMT}.
\end{equation}
Each arrow has an inverse on its image given by a projection of
coordinates. No output replaces its generative history.

\section{Theorem of internal gauge realization}

\begin{theorem}[Complete and natural internal realization]
\label{thm:gauint-realizacion}
For every \(x\in\mathcal C_{\rm cont}^{\rm gau}\), the chain \eqref{eq:gauint-cadena-completa} constructs, from the discrete structure of the continuum, a natural macrostructure with four directions, six incidences, eight faces, four involutions, a projective measure, a complete counting kernel, its positive transfer, a common eight-face law, a nonconstant incidence witness, and a cellular terminal. The thirteen coordinates of \eqref{eq:gauint-d-trece} remain recoverable in the output.
\end{theorem}

\begin{proof}
Equations \eqref{eq:gauint-cuatro-direcciones}--\eqref{eq:gauint-ocho-caras} derive \(4/6/8\). The incidence sum \eqref{eq:gauint-incidencia-invariante} proves the invariance of \(W_c\), and the uniform count proves \eqref{eq:gauint-wc-momentos}. The complete extensions produce \(k_m\) by \eqref{eq:gauint-kernel-normalizado}; the explicit domain conditions prove the Markov property, stationarity, and the two naturality equalities. Hence \(T_m=K_m^*K_m\) is positive and natural. Applying Fubini once to \eqref{eq:gauint-ley-ocho-caras} yields the four quadratic forms without changing the measure. The uniform marginal centers \(W_{c,m}\), whence \eqref{eq:gauint-w-eigen}. The initial object of the nerve gives the additional degeneracy and the SDR identity \eqref{eq:gauint-sdr}. Finally, the three graph constructions retain their inputs as first coordinates, so the composition forgets none of the thirteen inputs.
\end{proof}

\section{Provenance and Forgers}

\paragraph{Origin.}
The TPK seed, the oriented screen, the surviving extensions, the common measure, and the terminal contraction have the status \texttt{PREEXISTING\_AUTHORIAL\_ARCHITECTURE}. The counting transfer, the incidence mode, and the exact sector \(3Q\) are presented as a \texttt{RECOVERED\_RESULT}. The explicit derivation \(\mathbf P^1(\mathbf F_3)\to4/6/8\), the two lumpability conditions, and the non-ablative packaging by graphs have the status of a \texttt{NEW\_FORMALIZATION}; they formalize the preceding architecture without claiming a new conceptual origin.

\begin{ablation}[Internal falsifiers of the gauge branch]
\label{abl:gauint-falsadores}
The construction is falsified at the indicated level if any of the
following facts occurs.
\begin{enumerate}
 \item \(I_{\rm gau}\), \(E_{\rm gau}\), or
 \(\mathscr B_{\rm gau}\) is replaced by a list without the maps that derive them.
 \item \(sB=0\) fails or \(W_c\) ceases to be invariant and nonconstant.
 \item The count \(a_m(y,z)\) omits carry, orientation, memory, boundary, or
 route: in that case a different object has been counted.
 \item Projectivity, stationarity, or one of the two
 lumpabilities fails; then the corresponding naturality cannot be published.
 \item The eight faces do not follow from a unique \(\mu_m\) and a unique
 \(k_m\), or some \(\vartheta_L\) does not preserve \(\Omega_m^{(8)}\).
 \item \(P_{\Omega,m}W_{c,m}\ne0\) o
 \(D_m^{\rm HMT}W_{c,m}\ne3W_{c,m}\).
 \item The SDR identity \(\partial H+H\partial=I-\iota\epsilon\) fails.
 \item One of the thirteen coordinates is omitted: \emph{object replaced;
 conclusion not transferable}.
\end{enumerate}
These falsifiers control membership and naturality; they do not degrade a
history that satisfies all the identities.
\end{ablation}
\fi
